\documentclass[../../main.tex]{subfiles}
\begin{document}
% SOURCE p=240; printed=225; object=ch07-ex25-density-point-mass; number=25; region=lower
\begin{exercise}\label{exercise:ch07-density-point-mass}
设 $\mu$ 是 $X$ 上的 Radon 测度，且每个非空开集都有正测度（例如 Lebesgue 测度）。对每个 $x\in X$，存在 $L^1(\mu)$ 中的网 $\{f_\alpha\}$，在 $M(X)$ 中模糊收敛到 $x$ 处的点质量。若 $X$ 第一可数，可以取这个网为序列。（考虑形如 $\mu(U)^{-1}\chi_U$ 的函数。）
\end{exercise}
% SUPPLEMENT: id=v14-ch07-ex25-solution; parent=ch07-ex25-density-point-mass; kind=solution+motivation+hypotheses; origin=independent-derivation
\begin{solution}
\textbf{完整补写（重构新增）。}固定 $x\in X$。令 $\mathcal U$ 为所有包含 $x$ 的相对紧开邻域，按反包含排序：$U\preceq V$ 当且仅当 $V\subseteq U$。LCH 性保证这族邻域非空，并在每个开邻域中都能找到一个；两者的交仍是相对紧开邻域，所以它是有向集。每个 $U\in\mathcal U$ 满足 $0<\mu(U)\leqslant\mu(\overline U)<\infty$。令
\begin{align*}
f_U=\frac{\chi_U}{\mu(U)},\qquad\nu_U=f_U\mu.
\end{align*}
则 $f_U\in L^1(\mu)$，$\|f_U\|_1=1$，且由前面的嵌入 $\nu_U\in M(X)$。对任意 $g\in C_0(X)$ 和 $\varepsilon>0$，连续性给开邻域 $W\ni x$ 使 $|g(y)-g(x)|<\varepsilon$（$y\in W$）。选 $U_0\in\mathcal U$、$U_0\subseteq W$，则每个 $U\succeq U_0$ 都满足
\begin{align*}
\left|\int g\,\mathrm d\nu_U-g(x)\right|
\leqslant\frac1{\mu(U)}\int_U|g(y)-g(x)|\,\mathrm d\mu(y)
\leqslant\varepsilon.
\end{align*}
因此 $\nu_U\to\delta_x$ 模糊收敛。

若 $X$ 第一可数，在 $x$ 取可数开邻域基 $\{W_n\}$，再固定一个相对紧开邻域 $V\ni x$，令 $U_n=V\cap\bigcap_{j=1}^nW_j$。这些集合非空、开、相对紧且递减，仍构成 $x$ 的邻域基。相同估计说明 $\chi_{U_n}/\mu(U_n)$ 所对应的测度序列模糊趋于 $\delta_x$。
\end{solution}
\begin{remark}
\textbf{动机与条件（重构新增）。}把单位质量均匀放进越来越小的邻域，连续测试只看到邻域中心的值。每个非空开集正测度确保归一化分母不为零；紧集有限确保分母有限；局部紧性提供这样的邻域。第一可数只用于把邻域网换成序列。若某个包含 $x$ 的开集 $W$ 为零测集，取支撑于 $W$、在 $x$ 为 $1$ 的连续峰函数，则所有 $L^1(\mu)$ 密度的积分为零，而 $\delta_x$ 的读数为 $1$，所以此时不能得到所求逼近。
\end{remark}
% END SUPPLEMENT: id=v14-ch07-ex25-solution
% SOURCE p=240; printed=225; object=ch07-ex26-vague-norm-bounded-tests; number=26; region=lower
\begin{exercise}\label{exercise:ch07-vague-norm-bounded-tests}
若 $\{\mu_n\}\subseteq M(X)$，$\mu_n\to\mu$ 模糊收敛，且 $\|\mu_n\|\to\|\mu\|$，则对每个 $f\in BC(X)$，有 $\int f\,\mathrm d\mu_n\to\int f\,\mathrm d\mu$。（若 $\mu=0$，结论平凡；否则，对给定 $\varepsilon>0$，可取 $g\in C_c(X)$、$\|g\|_u\leqslant1$，使 $\int g\,\mathrm d\mu>\|\mu\|-\varepsilon$，且对 $f\in BC(X)$，$\int gf\,\mathrm d\mu_n\to\int gf\,\mathrm d\mu$。）此外，假设 $\|\mu_n\|\to\|\mu\|$ 不能省略。
\end{exercise}
% SUPPLEMENT: id=v14-ch07-ex26-solution; parent=ch07-ex26-vague-norm-bounded-tests; kind=solution+motivation+hypotheses+source-clarification; origin=independent-derivation
\begin{solution}
\textbf{完整补写（重构新增）。}若 $\mu=0$，则 $|\int f\,\mathrm d\mu_n|\leqslant\|f\|_u\|\mu_n\|\to0$。下面设 $\mu\ne0$。由 Riesz 等距定理及 $C_c$ 的一致稠密性，对任意 $\varepsilon>0$ 可取 $g\in C_c$、$\|g\|_u\leqslant1$，使
\begin{align*}
\operatorname{Re}\int g\,\mathrm d\mu>\|\mu\|-\varepsilon.
\end{align*}
确切地，可先在 $C_0$ 单位球内取近乎达到范数的测试函数，乘模为 $1$ 的常数使其积分为非负实数，再在 $C_c$ 中一致逼近并按其范数与 $1$ 的最大值缩放；误差可任意小。也可直接使用 Riesz 定理证明中的 Lusin 近似。原提示的实数不等号可按上述相位选择理解，或如这里明确写实部。

令 $K=\operatorname{supp}g$。因为 $g$ 在 $K$ 外为零，
\begin{align*}
\operatorname{Re}\int g\,\mathrm d\mu_n\leqslant|\mu_n|(K),\qquad
|\mu_n|(X\setminus K)\leqslant\|\mu_n\|-\operatorname{Re}\int g\,\mathrm d\mu_n.
\end{align*}
范数收敛与模糊收敛给
\begin{align*}
\limsup_n|\mu_n|(X\setminus K)
\leqslant\|\mu\|-\operatorname{Re}\int g\,\mathrm d\mu<\varepsilon.
\end{align*}
同样 $|\mu|(X\setminus K)<\varepsilon$。取 $\eta\in C_c(X)$ 满足 $0\leqslant\eta\leqslant1$、$\eta=1$ 于 $K$。任意 $f\in BC(X)$ 都有 $f\eta\in C_c(X)$，所以其积分按模糊收敛趋于极限，且
\begin{align*}
\left|\int f\,\mathrm d(\mu_n-\mu)\right|
&\leqslant\left|\int f\eta\,\mathrm d(\mu_n-\mu)\right|\\
&\quad+\|f\|_u\bigl(|\mu_n|(X\setminus K)+|\mu|(X\setminus K)\bigr).
\end{align*}
取上极限得到不超过 $2\varepsilon\|f\|_u$；令 $\varepsilon\downarrow0$ 即得所求。

反例取 $X=\mathbb R$、$\mu_n=\delta_n$、$\mu=0$。它们模糊收敛，但 $f\equiv1\in BC(\mathbb R)$ 的积分恒为 $1$，不趋于零；此时 $\|\mu_n\|=1\not\to0$。
\end{solution}
\begin{remark}
\textbf{动机与条件（重构新增）。}模糊收敛只测试在无穷远消失的函数；范数收敛排除有一部分总变差逃到每个紧集外的情形。上面的紧集 $K$ 使序列尾部与极限的外部质量同时小，才可把测试扩展到全部有界连续函数。有界性控制外部误差，连续性使截断后成为合法测试函数；一般有界可测测试仍不能替代，见\cref{exercise:ch07-vague-counterexamples}。
\end{remark}
% END SUPPLEMENT: id=v14-ch07-ex26-solution
% SOURCE p=240; printed=225; object=ch07-ex27-ck-dual; number=27; region=bottom
\begin{exercise}\label{exercise:ch07-ck-dual}
设 $C^k([0,1])$ 如\cref{exercise:ch05-ck-banach}。若 $I\in C^k([0,1])^*$，则存在 $\mu\in M([0,1])$ 和常数 $c_0,\ldots,c_{k-1}$，且它们全部唯一，使
\begin{align*}
I(f)=\int f^{(k)}\,\mathrm d\mu+\sum_{j=0}^{k-1}c_jf^{(j)}(0).
\end{align*}
（泛函 $f\mapsto f^{(j)}(0)$ 可换成任意一组在次数小于 $k$ 的多项式空间中分离点的 $k$ 个泛函。）
\end{exercise}
% SUPPLEMENT: id=v14-ch07-ex27-solution; parent=ch07-ex27-ck-dual; kind=solution+motivation+hypotheses+source-clarification; origin=independent-derivation
\begin{solution}
\textbf{完整补写（重构新增）。}范数沿用 $\|f\|_{C^k}=\sum_{j=0}^k\|f^{(j)}\|_u$。若 $k=0$，求和为空，结论就是紧区间的 Riesz 定理。以下 $k\geqslant1$。定义
\begin{align*}
Tf=(f^{(k)},f(0),f'(0),\ldots,f^{(k-1)}(0))
\in C([0,1])\oplus\mathbb C^k.
\end{align*}
右侧取范数 $\|(h,a)\|=\|h\|_u+\sum_{j=0}^{k-1}|a_j|$。显然 $\|Tf\|\leqslant\|f\|_{C^k}$。对任意 $(h,a)$，令
\begin{align*}
S(h,a)(x)=\sum_{j=0}^{k-1}a_j\frac{x^j}{j!}
+\frac1{(k-1)!}\int_0^x(x-t)^{k-1}h(t)\,\mathrm dt.
\end{align*}
逐次使用微积分基本定理，积分项的 $k$ 阶导数为 $h$，且其 $0$ 至 $k-1$ 阶导数在 $0$ 都为零，故 $TS=\operatorname{id}$。反复积分恢复函数的 Taylor 积分公式，又有 $ST=\operatorname{id}$。更具体地，对 $0\leqslant r<k$，
\begin{align*}
\|(S(h,a))^{(r)}\|_u
\leqslant\sum_{j=r}^{k-1}\frac{|a_j|}{(j-r)!}+\frac{\|h\|_u}{(k-r)!},
\end{align*}
而第 $k$ 阶范数等于 $\|h\|_u$。有限求和表明 $S$ 有界，所以 $T$ 是有界可逆坐标映射。

令 $J(h)=I(S(h,0))$，$c_j=I(x^j/j!)$。$S$ 的有界性保证 $J\in C([0,1])^*$，由 Riesz 定理存在唯一 $\mu\in M([0,1])$ 使 $J(h)=\int h\,\mathrm d\mu$。对 $f=S(Tf)$ 使用线性性，便得题示公式。若有另一组系数和测度，先对 $f=x^j/j!$ 测试，因 $f^{(k)}=0$ 且低阶原点导数为对应坐标，得到全部系数相同；再对所有 $f=S(h,0)$ 测试，Riesz 唯一性给测度相同。

最后补证括号中的推广，其中“泛函”按此对偶问题应取 $C^k([0,1])$ 上的\emph{连续线性}泛函 $\lambda_0,\ldots,\lambda_{k-1}$。设 $P_{k-1}$ 为次数小于 $k$ 的多项式空间。分离点意味着映射 $\Lambda:P_{k-1}\to\mathbb C^k$、$p\mapsto(\lambda_j(p))_j$ 单射；两边维数均为 $k$，所以它可逆，其逆在有限维空间上连续。令 $Rh=S(h,0)$。对任意 $h,a$，存在唯一多项式
\begin{align*}
p=\Lambda^{-1}\bigl(a-(\lambda_j(Rh))_j\bigr)
\end{align*}
使 $f=Rh+p$ 满足 $f^{(k)}=h$ 及 $(\lambda_j(f))_j=a$。由 $R$、各 $\lambda_j$ 及 $\Lambda^{-1}$ 的有界性，此逆构造有界；正向也有界。因此 $f\mapsto(f^{(k)},(\lambda_j(f))_j)$ 再次是 Banach 空间同构。重复前面的对偶分解和 Riesz 表示，就得到唯一的 $\mu,c_j$，使 $I(f)=\int f^{(k)}\,\mathrm d\mu+\sum_jc_j\lambda_j(f)$。
\end{solution}
\begin{remark}
\textbf{动机与逐条件说明（重构新增）。}最高阶导数决定函数至多相差一个低阶多项式；$k$ 个独立初始数据恰好消除这个核。$C^k$ 范数控制全部导数，使微分、初值与积分逆映射连续；只取一致范数做不到这一点。区间紧性允许对最高阶连续导数使用紧空间 Riesz 定理。推广中的分离条件保证低阶多项式可唯一恢复，而连续线性条件使新坐标仍是有界线性坐标；任意不连续或非线性的“读数”不能直接用于同一对偶表示。
\end{remark}
% END SUPPLEMENT: id=v14-ch07-ex27-solution
\end{document}
